\section{Abstract syntax trees and attribute grammars}
\label{sec:ast}

Some rules of a language cannot be captured by its syntax definition: type
consistency, or the use of \texttt{return} statements. Semantic analysis
checks them, and it works on a tree that is simpler than the parse tree. This
section and the next follow Wilhelm, Seidl and
Hack~\cite[Ch.~4]{wilhelm2013compiler} and Scott and
Aldrich~\cite[Ch.~4]{scott2025pragmatics}.

\subsection{Abstract syntax trees}

\begin{definition}
An \textbf{abstract syntax tree} (AST) is a parse tree with the nonessential
information removed: nonterminal nodes are replaced by nodes for the
constructs of the word they represent, and unit productions (rules such as
$\mathit{Exp} ::= \mathit{Term}$) are collapsed.
\end{definition}

Figure~\ref{fig:pt-vs-ast} compares the two trees for
\texttt{if x then y else z} in a grammar where the condition is an expression and
the branches are statements.
The parse tree has a node for every rule and every keyword; the AST has one
node for the conditional, with three children. ASTs are what compilers use to
generate code and what interpreters run.

\begin{figure}[h]
\centering
\begin{forest}
  for tree={font=\small\ttfamily, s sep=3mm, l sep=4mm},
  [S [if] [Exp [x]] [then] [S [y]] [else] [S [z]]]
\end{forest}
\hspace{2cm}
\begin{forest}
  for tree={font=\small\ttfamily, s sep=3mm, l sep=4mm},
  [if [x] [y] [z]]
\end{forest}
\caption{Parse tree (left) and AST (right) of a conditional.}
\label{fig:pt-vs-ast}
\end{figure}

\paragraph{Example.} For the imperative language of the semantics block, with
assignment $l := e$, sequence, conditional, \texttt{while} and dereferencing
$!l$, the program
\[
  \assign{y}{0}\;;\;
  \ift{\deref{x} = \deref{y}}
    {\assign{x}{\deref{x}+1}\;;\;\assign{y}{\deref{y}-1}}
    {\assign{x}{\deref{x}-1}\;;\;\assign{y}{\deref{y}+1}}
\]
has the AST of Figure~\ref{fig:simp-ast}. The top-level semicolon is the root.
$\deref{x}$ is a node \textsc{Deref} applied to the location $x$: the location names
a cell of the memory, and dereferencing reads its content.

\begin{figure}[h]
\centering
\begin{forest}
  for tree={font=\small, s sep=2mm, l sep=4mm},
  [\textsc{Seq}
    [\textsc{Assign} [$y$] [$0$]]
    [\textsc{If}
      [{$=$} [\textsc{Deref} [$x$]] [\textsc{Deref} [$y$]]]
      [\textsc{Seq}
        [\textsc{Assign} [$x$] [{$+$} [\textsc{Deref} [$x$]] [$1$]]]
        [\textsc{Assign} [$y$] [{$-$} [\textsc{Deref} [$y$]] [$1$]]]]
      [\textsc{Seq} [\dots]]
    ]
  ]
\end{forest}
\caption{AST of the example program. The \texttt{else} branch is symmetric to the \texttt{then} branch.}
\label{fig:simp-ast}
\end{figure}

\subsection{Attribute grammars}

\begin{definition}
An \textbf{attribute grammar} associates attributes with the symbols of a
grammar and, with each rule, semantic functions that compute the attributes
of the symbols of the rule. The attributes give the grammar elements a value
or a meaning.
\end{definition}

For the expression grammar, an attribute $\mathit{val}$ holds the value of each
expression. When a nonterminal appears twice in a rule, subscripts tell the
occurrences apart:
\begin{center}
\small
\begin{tabular}{l l}
  \toprule
  Rule & Semantic function \\
  \midrule
  $\mathit{Exp}_1 ::= \mathit{Exp}_2 + \mathit{Term}$ & $\mathit{Exp}_1.\mathit{val} := \mathit{sum}(\mathit{Exp}_2.\mathit{val}, \mathit{Term}.\mathit{val})$ \\
  $\mathit{Exp}_1 ::= \mathit{Exp}_2 - \mathit{Term}$ & $\mathit{Exp}_1.\mathit{val} := \mathit{difference}(\mathit{Exp}_2.\mathit{val}, \mathit{Term}.\mathit{val})$ \\
  $\mathit{Exp} ::= \mathit{Term}$ & $\mathit{Exp}.\mathit{val} := \mathit{Term}.\mathit{val}$ \\
  $\mathit{Term}_1 ::= \mathit{Term}_2 * \mathit{Factor}$ & $\mathit{Term}_1.\mathit{val} := \mathit{product}(\mathit{Term}_2.\mathit{val}, \mathit{Factor}.\mathit{val})$ \\
  $\mathit{Term}_1 ::= \mathit{Term}_2 / \mathit{Factor}$ & $\mathit{Term}_1.\mathit{val} := \mathit{quotient}(\mathit{Term}_2.\mathit{val}, \mathit{Factor}.\mathit{val})$ \\
  $\mathit{Term} ::= \mathit{Factor}$ & $\mathit{Term}.\mathit{val} := \mathit{Factor}.\mathit{val}$ \\
  $\mathit{Factor} ::= (\mathit{Exp})$ & $\mathit{Factor}.\mathit{val} := \mathit{Exp}.\mathit{val}$ \\
  $\mathit{Factor}_1 ::= - \mathit{Factor}_2$ & $\mathit{Factor}_1.\mathit{val} := \mathit{add\_inverse}(\mathit{Factor}_2.\mathit{val})$ \\
  $\mathit{Factor} ::= \mathit{Num}$ & $\mathit{Factor}.\mathit{val} := \mathit{Num}.\mathit{val}$ \\
  \bottomrule
\end{tabular}
\end{center}

\paragraph{Synthesized attributes.} Decorating (annotating) a parse tree starts
at the leaves and propagates up the tree following the semantic functions.
An attribute computed from the children of a node is \textbf{synthesized}.
Figure~\ref{fig:decorated} shows the decoration of $7 + 3 * 4$: the leaves
get the values of the numbers, the product node gets $12$, and the root gets
$7 + 12 = 19$.

\begin{figure}[h]
\centering
\begin{forest}
  for tree={font=\small, s sep=3mm, l sep=4mm},
  [{Exp \quad\textbf{19}}
    [{Exp \quad\textbf{7}} [{Term \quad\textbf{7}} [{Factor \quad\textbf{7}} [7]]]]
    [+]
    [{Term \quad\textbf{12}}
      [{Term \quad\textbf{3}} [{Factor \quad\textbf{3}} [3]]]
      [*]
      [{Factor \quad\textbf{4}} [4]]
    ]
  ]
\end{forest}
\caption{Parse tree of $7 + 3 * 4$ decorated with the synthesized attribute $\mathit{val}$.}
\label{fig:decorated}
\end{figure}

In Rascal, each semantic function is one equation over a concrete syntax
pattern, and the attribute is the function:
\codelines{Rascal}{rascal/src/Attributes.rsc}{22}{30}

\paragraph{Inherited attributes.} An attribute computed from the parent or the
siblings of a node is \textbf{inherited}. Inherited attributes carry values
from the left part of an expression into the right part, for example to
preserve associativity in a grammar without left recursion. With
\begin{center}
\small
\begin{tabular}{l l}
  \toprule
  Rule & Semantic functions \\
  \midrule
  $\mathit{Exp} ::= \mathit{Num}\ \mathit{ExpT}$ & $\mathit{ExpT}.\mathit{inh} := \mathit{Num}.\mathit{val}$;\ \ $\mathit{Exp}.\mathit{val} := \mathit{ExpT}.\mathit{syn}$ \\
  $\mathit{ExpT}_1 ::= - \mathit{Num}\ \mathit{ExpT}_2$ & $\mathit{ExpT}_2.\mathit{inh} := \mathit{ExpT}_1.\mathit{inh} - \mathit{Num}.\mathit{val}$;\ \ $\mathit{ExpT}_1.\mathit{syn} := \mathit{ExpT}_2.\mathit{syn}$ \\
  $\mathit{ExpT} ::= \lambda$ & $\mathit{ExpT}.\mathit{syn} := \mathit{ExpT}.\mathit{inh}$ \\
  \bottomrule
\end{tabular}
\end{center}
the value of $9 - 4 - 3$ is computed as in Figure~\ref{fig:inherited}. The 9
goes down into the first $\mathit{ExpT}$; there $9 - 4 = 5$ goes down into the
second; there $5 - 3 = 2$ goes down into the empty $\mathit{ExpT}$, which copies
its inherited value into its synthesized one; and the 2 goes back up to the
root. The subtraction associates to the left, $(9 - 4) - 3 = 2$, even though
the tree grows to the right. Computing $9 - (4 - 3) = 8$ would be the error of
ignoring associativity.

\begin{figure}[h]
\centering
\begin{forest}
  for tree={font=\small, s sep=4mm, l sep=5mm},
  [{Exp \quad syn 2}
    [{Num \quad 9}]
    [{ExpT \quad inh 9, syn 2}
      [$-$]
      [{Num \quad 4}]
      [{ExpT \quad inh 5, syn 2}
        [$-$]
        [{Num \quad 3}]
        [{ExpT \quad inh 2, syn 2} [$\lambda$]]
      ]
    ]
  ]
\end{forest}
\caption{Decoration of $9 - 4 - 3$ with an inherited attribute ($\mathit{inh}$) and a synthesized one ($\mathit{syn}$).}
\label{fig:inherited}
\end{figure}

In Rascal the inherited attribute is an extra parameter:
\codelines{Rascal}{rascal/src/Attributes.rsc}{35}{38}

\subsection{Building ASTs}

Attribute grammars can also build the AST from the parse tree. An attribute
$\mathit{ptr}$ points to the AST node of each grammar symbol, and three
constructors create nodes: $\mathit{make\_leaf}(n)$, $\mathit{make\_bin\_op}(op, l, r)$
and $\mathit{make\_un\_op}(op, a)$.

\paragraph{Bottom-up.} With the grammar with left recursion, $\mathit{ptr}$ is
synthesized:
\begin{center}
\small
\begin{tabular}{l l}
  \toprule
  Rule & Semantic function \\
  \midrule
  $\mathit{Exp}_1 ::= \mathit{Exp}_2 + \mathit{Term}$ & $\mathit{Exp}_1.\mathit{ptr} := \mathit{make\_bin\_op}(``+", \mathit{Exp}_2.\mathit{ptr}, \mathit{Term}.\mathit{ptr})$ \\
  $\mathit{Exp}_1 ::= \mathit{Exp}_2 - \mathit{Term}$ & $\mathit{Exp}_1.\mathit{ptr} := \mathit{make\_bin\_op}(``-", \mathit{Exp}_2.\mathit{ptr}, \mathit{Term}.\mathit{ptr})$ \\
  $\mathit{Exp} ::= \mathit{Term}$ & $\mathit{Exp}.\mathit{ptr} := \mathit{Term}.\mathit{ptr}$ \\
  $\mathit{Term}_1 ::= \mathit{Term}_2 * \mathit{Fact}$ & $\mathit{Term}_1.\mathit{ptr} := \mathit{make\_bin\_op}(``*", \mathit{Term}_2.\mathit{ptr}, \mathit{Fact}.\mathit{ptr})$ \\
  $\mathit{Term}_1 ::= \mathit{Term}_2 / \mathit{Fact}$ & $\mathit{Term}_1.\mathit{ptr} := \mathit{make\_bin\_op}(``/", \mathit{Term}_2.\mathit{ptr}, \mathit{Fact}.\mathit{ptr})$ \\
  $\mathit{Term} ::= \mathit{Fact}$ & $\mathit{Term}.\mathit{ptr} := \mathit{Fact}.\mathit{ptr}$ \\
  $\mathit{Fact} ::= (\mathit{Exp})$ & $\mathit{Fact}.\mathit{ptr} := \mathit{Exp}.\mathit{ptr}$ \\
  $\mathit{Fact}_1 ::= - \mathit{Fact}_2$ & $\mathit{Fact}_1.\mathit{ptr} := \mathit{make\_un\_op}(``-", \mathit{Fact}_2.\mathit{ptr})$ \\
  $\mathit{Fact} ::= \mathit{Num}$ & $\mathit{Fact}.\mathit{ptr} := \mathit{make\_leaf}(\mathit{Num}.\mathit{val})$ \\
  \bottomrule
\end{tabular}
\end{center}
The slide assigns $\mathit{Exp}_1.\mathit{val}$ in the rule for subtraction; the
attribute is $\mathit{ptr}$, since the rule builds a node. It also writes the
unary rule with $\mathit{Factor}$ on the left and $\mathit{Fact}$ on the right; the
table uses $\mathit{Fact}$ throughout.

For $(1+2)*4$ the leaves become nodes first, then the sum, then the product
(Figure~\ref{fig:ast-bu}). The parentheses and the unit rules disappear: they
only pass pointers up.

\begin{figure}[h]
\centering
\begin{forest}
  for tree={font=\small, s sep=3mm, l sep=4mm},
  [Exp
    [Term
      [Term [Fact [( ] [Exp [Exp [Term [Fact [Num [1]]]]] [+] [Term [Fact [Num [2]]]]] [ )]]]
      [*]
      [Fact [Num [4]]]
    ]
  ]
\end{forest}
\hspace{1.5cm}
\begin{forest}
  for tree={font=\small\ttfamily, s sep=4mm, l sep=6mm},
  [$*$ [$+$ [1] [2]] [4]]
\end{forest}
\caption{Parse tree of $(1+2)*4$ and the AST its $\mathit{ptr}$ attributes build. The $+$ node is created when the parser reduces $\mathit{Exp} ::= \mathit{Exp} + \mathit{Term}$ inside the parentheses, the $*$ node at $\mathit{Term} ::= \mathit{Term} * \mathit{Fact}$.}
\label{fig:ast-bu}
\end{figure}

\codelines{Rascal}{rascal/src/Attributes.rsc}{42}{46}
\codelines{Rascal}{rascal/src/Attributes.rsc}{48}{56}

\paragraph{Top-down.} With the grammar without left recursion the operator and
its left operand are in different subtrees, so the tree built so far has to
travel down to the right as an inherited attribute $\mathit{st}$:
\begin{center}
\small
\begin{tabular}{l >{\raggedright\arraybackslash}p{0.6\textwidth}}
  \toprule
  Rule & Semantic functions \\
  \midrule
  $\mathit{Exp} ::= \mathit{Term}\ \mathit{Terms}$ & $\mathit{Terms}.\mathit{st} := \mathit{Term}.\mathit{ptr}$;\newline $\mathit{Exp}.\mathit{ptr} := \mathit{Terms}.\mathit{ptr}$ \\
  $\mathit{Terms}_1 ::= + \mathit{Term}\ \mathit{Terms}_2$ & $\mathit{Terms}_2.\mathit{st} := \mathit{make\_bin\_op}(``+", \mathit{Terms}_1.\mathit{st}, \mathit{Term}.\mathit{ptr})$;\newline $\mathit{Terms}_1.\mathit{ptr} := \mathit{Terms}_2.\mathit{ptr}$ \\
  $\mathit{Terms}_1 ::= - \mathit{Term}\ \mathit{Terms}_2$ & same, with $``-"$ \\
  $\mathit{Terms} ::= \lambda$ & $\mathit{Terms}.\mathit{ptr} := \mathit{Terms}.\mathit{st}$ \\
  $\mathit{Term} ::= \mathit{Fact}\ \mathit{FactTerm}$ & $\mathit{FactTerm}.\mathit{st} := \mathit{Fact}.\mathit{ptr}$;\newline $\mathit{Term}.\mathit{ptr} := \mathit{FactTerm}.\mathit{ptr}$ \\
  $\mathit{FactTerm}_1 ::= * \mathit{Fact}\ \mathit{FactTerm}_2$ & $\mathit{FactTerm}_2.\mathit{st} := \mathit{make\_bin\_op}(``*", \mathit{FactTerm}_1.\mathit{st}, \mathit{Fact}.\mathit{ptr})$;\newline $\mathit{FactTerm}_1.\mathit{ptr} := \mathit{FactTerm}_2.\mathit{ptr}$ \\
  $\mathit{FactTerm}_1 ::= / \mathit{Fact}\ \mathit{FactTerm}_2$ & same, with $``/"$ \\
  $\mathit{FactTerm} ::= \lambda$ & $\mathit{FactTerm}.\mathit{ptr} := \mathit{FactTerm}.\mathit{st}$ \\
  $\mathit{Fact}_1 ::= - \mathit{Fact}_2$ & $\mathit{Fact}_1.\mathit{ptr} := \mathit{make\_un\_op}(``-", \mathit{Fact}_2.\mathit{ptr})$ \\
  $\mathit{Fact} ::= (\mathit{Exp})$ & $\mathit{Fact}.\mathit{ptr} := \mathit{Exp}.\mathit{ptr}$ \\
  $\mathit{Fact} ::= \mathit{Num}$ & $\mathit{Fact}.\mathit{ptr} := \mathit{make\_leaf}(\mathit{Num}.\mathit{val})$ \\
  \bottomrule
\end{tabular}
\end{center}
Four rules on the slide differ from this table. In the first rule it assigns
$\mathit{Terms}.\mathit{ptr} := \mathit{Term}.\mathit{ptr}$, but the left operand has to go
into the inherited $\mathit{st}$, as in the other rules of $\mathit{Terms}$. The same
happens with $\mathit{FactTerm}.\mathit{st}$ (written $\mathit{FafctTerm}.\mathit{ptr}$). The rules
for $*$ and $/$ compute $\mathit{FactTerm}_1.\mathit{st}$ from $\mathit{Fact}.\mathit{st}$ and
$\mathit{FactTerm}_2.\mathit{ptr}$, the reverse of the rules for $+$ and $-$; the node
must be built from the inherited tree and the new factor and passed down to
$\mathit{FactTerm}_2$. And the parenthesis rule is written $\mathit{Fact} ::= \mathit{Exp}$.
\codelines{Rascal}{rascal/src/Attributes.rsc}{58}{71}
The tests check that both constructions build the same AST for $(1+2)*4$ and
for $1-2*3+4/2$, and that $9-4-3$ becomes $(9-4)-3$.

\subsection{Exercises}

\begin{exercise}
Decorate the parse tree of $8 / 2 - (1 + 1)$ with the synthesized attribute
$\mathit{val}$.
\end{exercise}

\paragraph{Solution.} The leaves get 8, 2, 1 and 1. $\mathit{Term} ::= \mathit{Term} / \mathit{Factor}$
gives $8 / 2 = 4$. Inside the parentheses, $\mathit{Exp} ::= \mathit{Exp} + \mathit{Term}$ gives
$2$, and $\mathit{Factor} ::= (\mathit{Exp})$ copies it. The root,
$\mathit{Exp} ::= \mathit{Exp} - \mathit{Term}$, gives $4 - 2 = 2$.

\begin{exercise}
Extend the inherited-attribute grammar for $9-4-3$ with addition, and compute
$10 - 3 + 2$.
\end{exercise}

\paragraph{Solution.} Add the rule $\mathit{ExpT}_1 ::= + \mathit{Num}\ \mathit{ExpT}_2$ with
\[
\mathit{ExpT}_2.\mathit{inh} := \mathit{ExpT}_1.\mathit{inh} + \mathit{Num}.\mathit{val}
\qquad
\mathit{ExpT}_1.\mathit{syn} := \mathit{ExpT}_2.\mathit{syn}.
\]
The inherited values are 10, then
$10 - 3 = 7$, then $7 + 2 = 9$, and 9 goes up to the root.

\begin{exercise}
Build the AST of $\whiledo{\deref{i}<\deref{n}}{\seq{\assign{s}{\deref{s}+\deref{i}}}{\assign{i}{\deref{i}+1}}}$.
\end{exercise}

\paragraph{Solution.} The root is \textsc{While} with two children: the
condition $<(\textsc{Deref}(i), \textsc{Deref}(n))$ and the body
$\textsc{Seq}(\textsc{Assign}(s, +(\textsc{Deref}(s), \textsc{Deref}(i))),
\textsc{Assign}(i, +(\textsc{Deref}(i), 1)))$.

\begin{exercise}
Run the top-down construction on $6 / 3 / 2$ and show the value of
$\mathit{st}$ in each $\mathit{FactTerm}$.
\end{exercise}

\paragraph{Solution.} The first $\mathit{FactTerm}$ inherits the leaf 6. It reads
$/ 3$ and passes $\mathit{st} = /(6, 3)$ to the second, which reads $/ 2$ and
passes $/(/(6, 3), 2)$ to the empty $\mathit{FactTerm}$. That tree goes back up as
$\mathit{ptr}$: the division associates to the left, $(6/3)/2 = 1$.
